\subsection{A full-mass point pole without a smooth maximizing test}
\label{subsec:point-pole-dual-nonattainment}

The smaller-background criteria do not require $L^p$ integrability. The following closed-background example has genuine obstacle envelopes at every scale and an unbounded threshold solution. Its dual supremum is approached by smooth tests but is not attained by any smooth semipositive test.

\begin{proposition}\label{prop:point-pole-full-mass-nonattainment}\index{Gauduchon duality!point-pole nonattainment}\index{Monge--Amp\`ere density!point-pole integrability}
Let $X=\mathbb P^n$, $n\geq2$, and let $\omega$ be the normalized Fubini--Study form, with $\int_X\omega^n=1$ and $\mathrm{d}V=\omega^n$. For $0<\alpha<1$ put
\[
 \psi=\log\frac{\sum_{i=1}^n|z_i|^2}{\sum_{i=0}^n|z_i|^2},
 \qquad u=1-(1-\psi)^\alpha,
 \qquad o=[1:0:\cdots:0].
\]
The function $u$ is a genuine normalized $\omega$-plurisubharmonic potential with sole pole $o$ and extended mass one. Every $P_{\delta\omega}(u)$ is genuine, for $0<\delta<1$. Its measure is $f\,\mathrm{d}V$ with $f\in L^1(\mathrm{d}V)$ and $f\notin L^p(\mathrm{d}V)$ for every $p>1$.

One has $b_\omega(f)=1$ and $\rho_\omega(f^{1/n})=1$. Thus $u$ is an actual threshold solution. No smooth semipositive $T\in\mathscr G_\omega$ attains $D_{f^{1/n}}(T)=1$.
\end{proposition}
\begin{proof}
The numerator vanishes only at $o$. Locally $\omega+\ddc\psi=\ddc\log\sum_{i=1}^n|z_i|^2\geq0$. The function $\chi(t)=1-(1-t)^\alpha$ is increasing and convex on $t\leq0$, with $0\leq\chi'\leq\alpha$. Apply its chain inequality to bounded truncations of $\psi$ and take the decreasing limit. It gives a genuine $u\in\PSH(X,\omega)$, with
\[
 A=\omega+\ddc u\geq(1-\alpha)\omega.
\]
The function is smooth off $o$ and has supremum zero on $\{z_0=0\}$. Since $u=-(1-\psi)^\alpha+1$, \cref{prop:sublinear-threshold-envelope} gives all the smaller-background envelopes.

To compute its mass, use the chart $z_0=1$ and write $s=\sum_{i=1}^n|z_i|^2$. A local potential for $A$ and its radial flux are
\begin{equation}\label{eq:point-pole-radial-flux}
 \begin{aligned}
 F(s)&=\log(1+s)+1-\left(1-\log\frac{s}{1+s}\right)^\alpha,\\
 q(s)&=sF'(s)=\frac{s+\alpha(1-\log(s/(1+s)))^{\alpha-1}}{1+s}.
 \end{aligned}
\end{equation}
The tangential Hessian eigenvalue is $q/s$ and the radial one is $q'$. Polar integration gives the radial top-degree measure $nq^{n-1}q'\,\mathrm{d}s$. Its angular normalization follows from the Fubini--Study case $q=s/(1+s)$, whose radial density is $ns^{n-1}(1+s)^{-n-1}\,\mathrm{d}s$ and whose total mass is one. Consequently
\[
 \int_{\{a<|z|^2<b\}}A^n=q(b)^n-q(a)^n.
\]
Here $q(s)\to0$ as $s\downarrow0$ and $q(s)\to1$ as $s\to\infty$. On each compact annulus the extended measure equals the smooth top power. It gives no mass to $o$ or to the infinity hyperplane, since both are pluripolar. Exhausting the affine chart by annuli therefore gives extended mass exactly one. This is a mass calculation, independent of the vanishing of Lelong numbers.

Near $o$ write $L=1-\log s$. Direct differentiation of \eqref{eq:point-pole-radial-flux} gives
\[
 \begin{gathered}
 \frac qs\sim\alpha s^{-1}L^{\alpha-1},\qquad
 q'\sim\alpha(1-\alpha)s^{-1}L^{\alpha-2},\\
 f(s)\asymp s^{-n}L^{-1-n(1-\alpha)}.
 \end{gathered}
\]
The smooth radial volume is comparable to $s^{n-1}\,\mathrm{d}s$. Its $L^1$ tail is thus comparable to $\int^\infty L^{-1-n(1-\alpha)}\,\mathrm{d}L$, which is finite. For every $p>1$, the corresponding tail has an exponentially growing factor $\exp(n(p-1)L)$ and diverges. The measure is smooth off $o$, so it is globally $f\,\mathrm{d}V$. Since its total mass is one, \cref{thm:invariant-mass-threshold-envelopes} gives $b_\omega(f)=1$, and \cref{thm:integrable-gauduchon-formula} gives $\rho_\omega(f^{1/n})=1$.

Suppose a smooth semipositive normalized test $T$ attained this value. The positive current $A$ satisfies $\int_XA\wedge T=1$ by distributional Stokes. On $X\setminus\{o\}$, matrix arithmetic--geometric mean gives
\[
 A\wedge T\geq f^{1/n}k_T\,\mathrm{d}V.
\]
The right side has total mass one by attainment. Thus the nonnegative gap vanishes off $o$. There $A$ is smooth strictly positive, so equality gives $T=tA^{n-1}$, including at points where $T=0$. This equality case follows by applying arithmetic--geometric mean to the positive semidefinite coefficient matrix conjugated by $A^{1/2}$: all its eigenvalues agree, possibly at zero. The function $t=(T\wedge A)/A^n$ is smooth and nonnegative on the punctured manifold.

Since $A$ is closed, $\ddc T=0$ becomes $\ddc t\wedge A^{n-1}=0$, a locally strictly elliptic scalar equation there. All eigenvalues of $A$ tend to infinity at $o$, so, as $n\geq2$, all complementary-product eigenvalues of $A^{n-1}$ do also. Boundedness of the smooth coefficient matrix of $T$ forces $t\to0$ at $o$.

Extend $t(o)=0$. This continuous nonnegative function on compact $X$ is nonzero by the normalization of $T$. It has a positive maximum away from $o$. The local strong maximum principle on connected $X\setminus\{o\}$ forces it to be constant there, contrary to its zero limit at $o$. This excludes every smooth semipositive maximizing test; no assumption about an analytic degeneracy set was needed.
\end{proof}

\begin{corollary}\label{cor:point-pole-entropy-energy-threshold}\index{Entropy!point-pole threshold}
For the density and potential of \cref{prop:point-pole-full-mass-nonattainment},
\[
 \int_X f\log f\,\mathrm{d}V<+\infty
 \quad\Longleftrightarrow\quad \alpha<1-\frac1n.
\]
For every real $r>0$,
\[
 \int_X(-u)^r f\,\mathrm{d}V<+\infty
 \quad\Longleftrightarrow\quad \alpha r<n(1-\alpha).
\]
In particular, finite entropy of this density can coexist with an unbounded full-mass threshold solution and absence of a smooth semipositive maximizing test.
\end{corollary}
\begin{proof}
All possible divergence is near $o$. The preceding estimates give $f\,\mathrm{d}V\asymp L^{-1-n(1-\alpha)}\,\mathrm{d}L$, $\log f\sim nL$, and $-u\asymp L^\alpha$. The entropy tail converges exactly when $n(1-\alpha)>1$; the weighted tail converges exactly when $\alpha r<n(1-\alpha)$. At equality each tail is comparable to $\int^\infty L^{-1}\,\mathrm{d}L$ and diverges. The negative part of $f\log f$ is boundedly integrable on the probability volume, so the entropy criterion is unambiguous.
\end{proof}

These are entropy and weighted integrals of the single equation density. The original S16 question concerns uniform moment or entropy conditions for determinant-weighted probabilities of a maximizing sequence under smooth semipositive attainment. The present example lies outside that attained setting. The proposition concerns the explicit threshold solution $u$; its identification with the greatest threshold envelope would require a separate uniqueness argument.
